% year: תשפ"ד
% type: מבחן
% semester: א
% moed: ב
% number: 2
% points: 16
% categories: continuity, func-limits
% source: exams/מבחנים/תשפד/מועד ב תשפד.pdf

\begin{question}
נגדיר פונקציה על ידי
\[ f(x)=\begin{cases}\frac{\ln(1+x^2)}{\sin^2 x}, & x<0,\\ 0, & 0\le x\le 1,\\ \frac{x}{x^2-1}, & x>1\end{cases} \]
מצאו את כל נקודות אי הרציפות של הפונקציה, ומיינו אותן לאי רציפות מסוג סליקה, קפיצה או עיקרית. נמקו.
\end{question}

\begin{hint}
בתוך כל אחד מהתחומים הפונקציה היא הרכבה/מנה של פונקציות אלמנטריות. בדקו היכן מתאפס מכנה, ובנפרד את נקודות התפר $x=0$ ו-$x=1$.
\end{hint}

\begin{hint}
ליד $0$: $\ln(1+x^2)\approx x^2$ ו-$\sin^2x\approx x^2$.
\end{hint}

\begin{solution}
\textbf{בתוך התחומים הפתוחים.}
\begin{itemize}
  \item בקטע $(0,1)$ הפונקציה קבועה ולכן רציפה.
  \item בקטע $(1,\infty)$: $f(x)=\frac{x}{x^2-1}$ היא מנה של פולינומים שהמכנה שלה אינו מתאפס שם ($x^2-1>0$), ולכן רציפה.
  \item בקטע $(-\infty,0)$: $f(x)=\frac{\ln(1+x^2)}{\sin^2x}$ היא מנה של פונקציות רציפות, ולכן רציפה בכל נקודה שבה $\sin x\neq0$. המכנה מתאפס בנקודות $x=-k\pi$, $k=1,2,3,\dots$, ושם הפונקציה אינה מוגדרת. בנקודות אלה המונה שואף ל-$\ln(1+k^2\pi^2)>0$ והמכנה שואף ל-$0^+$ (כי $\sin^2x>0$ בסביבה מנוקבת), ולכן
  \[ \lim_{x\to-k\pi^\pm}f(x)=+\infty. \]
  הגבולות החד-צדדיים אינם סופיים, ולכן $x=-k\pi$ ($k\in\N$) הן נקודות אי-רציפות \textbf{עיקריות} (מסוג שני).
\end{itemize}

\textbf{הנקודה $x=0$.} $f(0)=0$, ומימין $f\equiv0$ ולכן $\lim_{x\to0^+}f(x)=0$. משמאל:
\[ \lim_{x\to0^-}\frac{\ln(1+x^2)}{\sin^2x}=\lim_{x\to0^-}\frac{\ln(1+x^2)}{x^2}\cdot\left(\frac{x}{\sin x}\right)^2=1\cdot1^2=1, \]
כאשר השתמשנו בגבול היסודי $\lim_{t\to0}\frac{\ln(1+t)}{t}=1$ (עם $t=x^2\to0$) וב-$\lim_{x\to0}\frac{\sin x}{x}=1$. שני הגבולות החד-צדדיים סופיים ושונים ($1\neq0$), ולכן ב-$x=0$ יש אי-רציפות מסוג \textbf{קפיצה}.

\textbf{הנקודה $x=1$.} $f(1)=0$ ו-$\lim_{x\to1^-}f(x)=0$. מימין, המונה שואף ל-$1$ והמכנה $x^2-1\to0^+$, לכן
\[ \lim_{x\to1^+}\frac{x}{x^2-1}=+\infty. \]
הגבול החד-צדדי מימין אינו סופי, ולכן ב-$x=1$ יש אי-רציפות \textbf{עיקרית}.

\textbf{תשובה:} ב-$x=0$ -- אי-רציפות מסוג קפיצה; ב-$x=1$ -- אי-רציפות עיקרית; בנקודות $x=-k\pi$, $k=1,2,\dots$ (שבהן $f$ אינה מוגדרת) -- אי-רציפות עיקרית (הגבול שם הוא $+\infty$). בכל שאר הנקודות $f$ רציפה.
\end{solution}
